\metatitle{Interlacing polynomials}

\metadescription{Theory of interlacing and interleaving roots, Sturm sequences, compatible polynomials, and operators preserving real-rootedness.}


\section[interlacingPolynomials]{Interlacing and interleaving roots}


The following definitions are from \cite{Wagner1992}.  See also
\cite{Branden2015}.

\begin{definition}[Interleaving polynomials]
Let $f$ and $g$ be polynomials with positive leading coefficients and
with real roots $\{f_{i}\}$ and $\{g_{i}\}$, respectively.
We say that $f$ \defin{interlaces} $g$ if $\deg(f)+1=\deg(g)=d$ and
\[
g_1 \leq f_1 \leq g_2 \leq f_2 \leq \dotsm \leq  f_{d-1}\leq g_{d}.
\]
Moreover, we say that $f$ \defin{alternates left of} $g$ if $\deg(f)=\deg(g)=d$ and
\[
f_1 \leq g_1 \leq f_2 \leq \dotsm \leq f_{d}\leq g_{d}.
\]
We say that $f$ \defin{interleaves} $g$ if either $f$ interlaces $g$ or
$f$ alternates left of $g$. We write this as $f \interl g$.
By convention, $0 \interl 0$, $0 \interl h$ and $h \interl 0$ whenever $h$ is a polynomial with
a positive leading coefficient.
\end{definition}


\begin{lemma}[Folklore]
    Let $f, g$ be monic polynomials with $\deg(f) = \deg(g)+1$ and all real, non-positive roots.
    If $g \interl f$, then $f-tg$ is real-rooted and $g \interl f- tg \interl f$.
\end{lemma}


The following very useful lemma is from \cite[Sec. 3]{Wagner1992}.
\begin{lemma}[Wagner's lemma]
Let $f, g, h \in \setR[t]$ be real-rooted polynomials with
only real, non-positive roots and positive leading coefficients.
Then
\begin{itemize}
\item if $f \interl h$ and $g \interl h$ then $f+g \interl h$.
\item if $h \interl f$ and $h \interl g$ then $h \interl f+g$.
\item $g \interl f$ if and only if $f \interl tg$.
\end{itemize}
\end{lemma}

\name[D.G. Wagner]{David Wagner} also shows that $f \interl g$ implies
that $f \interl \lambda f + \mu g \interl g$ for all $\lambda, \mu \geq 0$.


As an example, if $f$ is real-rooted, then $f' \interl f$.
Hence, recursions involving the derivative play well with interlacing roots.


\begin{theorem}[Cauchy interlacing theorem, see \cite{Godsil2017,Fisk2005x}]
Let $A$ be a real symmetric $n \times n$ matrix with eigenvalues
\[
 \lambda_1 \leq \lambda_2 \leq \dotsb \leq \lambda_n,
\]
and let $B$ be a principal $(n-1)\times(n-1)$ submatrix of $A$ with eigenvalues
\[
 \mu_1 \leq \mu_2 \leq \dotsb \leq \mu_{n-1}.
\]
Then
\[
 \lambda_1 \leq \mu_1 \leq \lambda_2 \leq \mu_2 \leq \dotsb
 \leq \mu_{n-1} \leq \lambda_n.
\]
Equivalently, the characteristic polynomial of $B$ interlaces the
characteristic polynomial of $A$.
\end{theorem}



A more recent result is the following theorem by \name{Shi-Mei Ma} and \name{Yi Wang} \cite{MaWang2008}.
\begin{theorem}[Ma and Wang, 2008]
Suppose $f$ and $F$ are polynomials with positive leading coefficients,
such that
\[
 F(x) = u(x) f(x) + v(x) f'(x), \qquad  \deg(f) \leq \deg(F) \leq \deg(f) + 1
 \]
with $u(x),v(x)$ having real coefficients.
If $f$ is real-rooted and $v(r) \neq 0$ whenever $f(r)=0$,
then $F$ is real-rooted and $f \interl F$.
\end{theorem}
The theorem states more regarding root placements and multiplicities.


The following criterion from \cite[Thm. 2.1]{LiuWang2007} is often useful.
\begin{theorem}[Liu--Wang, 2007]
Let $F,f,g \in \setR[t]$ satisfy
\[
 F(x) = a(x)f(x) + b(x)g(x), \qquad \deg(F) \in \{\deg(f), \deg(f)+1\},
\]
where $a(x),b(x) \in \setR[t]$.
Assume that $f$ and $g$ are real-rooted, that $g \interl f$, and that $F$ and $g$ have leading coefficients of the same sign.
If $b(r) \leq 0$ whenever $f(r)=0$, then $F$ is real-rooted and $f \interl F$.
In particular, if $g$ strictly interlaces $f$ and $b(r) \lt 0$ whenever $f(r)=0$, then $f$ strictly interlaces $F$.
\end{theorem}



\subsection[sturmSequence]{Sturm sequences}


It is fairly common to see sequences of polynomials, $\{ P_n(t) \}_{n \geq 0}$
with the property that $P_j(t) \interl P_{j+1}(t)$ for all $j$.
If the leading coefficient is always positive and $\deg(P_n)=n$,
this is called a \defin{Sturm sequence}, see \cite{LiuWang2007}.
If we drop the degree condition, the sequence is called a \defin{generalized Sturm sequence}.

The notion of \defin{Sturm-unimodality} is defined in the analogous manner.



There are many Sturm sequences on \hyperref[realRootedWords]{permutations} and \hyperref[realRootedCatalan]{Catalan objects}.


\subsection[bezoutMatrix]{Bézout matrices and Hurwitz--Lace}

For arbitrary polynomials $f(t)$ and $g(t)$ with $\deg(f) = \deg(g)+1$,
let $B(f,g)$ be the \defin{Bézout matrix}, whose $(i,j)$ entry is equal to the coefficient of $x^{i-1} y^{j-1}$ in the polynomial
\[
 \frac{f(x)g(y) - f(y)g(x)}{x-y}.
\]

\begin{theorem}[Bézout's theorem, {\cite[Cor. 9.145]{Fisk2006x}}]
Let $f(t)$ be a polynomial of degree $d$ and $g(t)$ a polynomial of degree $d-1$.
Then $f(t)$ and $g(t)$ are both real-rooted and $g$ strictly interlaces $f$
if and only if $B(f,g)$ is positive definite.
\end{theorem}



\begin{theorem}[Hurwitz--Lace criterion]
Let
\[
p_0(t)=a_0+a_1t+a_2t^2+\dotsb,\qquad
p_1(t)=b_0+b_1t+b_2t^2+\dotsb
\]
have non-negative coefficients. Form the Hurwitz matrix
\[
H(p_0,p_1)=
\begin{pmatrix}
b_0 & b_1 & b_2 & b_3 & \dotsb \\
a_0 & a_1 & a_2 & a_3 & \dotsb \\
0   & b_0 & b_1 & b_2 & \dotsb \\
0   & a_0 & a_1 & a_2 & \dotsb \\
0   & 0   & b_0 & b_1 & \dotsb \\
0   & 0   & a_0 & a_1 & \dotsb \\
\vdots & \vdots & \vdots & \vdots & \ddots
\end{pmatrix}.
\]
Then \(H(p_0,p_1)\) is \hyperref[totallyNonnegative]{totally non-negative} if and only if \(p_1 \interl p_0\).
\end{theorem}

This criterion packages interlacing roots as total non-negativity of a
structured Hurwitz matrix.

\begin{theorem}[Hermite--Biehler, see \cite{Holtz2003}]
Let
\[
  P(z)=a_0+a_1z+\dotsb+a_nz^n
\]
be a real polynomial, and define its even and odd parts by
\[
  P_E(t)=\sum_{k\geq 0} a_{2k}t^k,
  \qquad
  P_O(t)=\sum_{k\geq 0} a_{2k+1}t^k.
\]
Assume that $P_E P_O\not\equiv 0$ and that $P_E$ and $P_O$ have
positive leading coefficients. Then $P$ has all zeros in the closed
left half-plane if and only if $P_E$ and $P_O$ have only real
non-positive zeros and $P_O \interl P_E$.
\end{theorem}

Thus the \name{Hermite--Biehler} and \name{Routh--Hurwitz} criteria
also translate root-location questions into interlacing and total
non-negativity.  The block-Toeplitz and network viewpoint is discussed
in \cite[Sec.~4]{Angarone2023x}, and weak/quasi-stable boundary cases
are treated in \cite{AdmGarloffTyaglov2017}.



\section[interlacingSequences]{Interlacing sequences}


The following definitions are from \cite{Branden2015}.

\begin{definition}[Interlacing sequences of polynomials]
A sequence $\{f_1,f_2,\dotsc,f_n\}$ of polynomials with positive leading coefficients
is said to be an \defin{interlacing sequence}
if $f_i \interl f_j$  for all $1\leq i \lt j \leq n$.

We let \defin{$\mathscr{F}_n^+$} denote the family of interlacing
sequences of length $n$ with \emph{non-negative} coefficients.
\end{definition}

In \cite{Branden2015}, there are nice examples of how interlacing
sequences can be used to prove real-rootedness.
In particular, \hyperref[realRootedMatching]{rook polynomials} associated with
Ferrers diagrams can be proved to be real-rooted using this method.


\begin{proposition}[\cite{Wagner1992}]
Suppose $f_1 \interl f_2 \interl \dotsb \interl f_n$ and $f_1 \interl f_n$.
Then $\{f_1,\dotsc,f_n\}$ is an interlacing sequence.
\end{proposition}


\subsection[fullyInterlacingLace]{Fully interlacing sequences and the $\operatorname{Lace}$ matrix}

A stronger coefficient-level version of interlacing was introduced by
\name{Christos A. Athanasiadis} and \name{David G. Wagner}
\cite{AthanasiadisWagner2024x}.

\begin{definition}[The $\operatorname{Lace}$ matrix]
Let $A=(A_{ij}(x))$ be a $p \times q$ matrix of formal power series, indexed by
$0 \leq i \lt p$ and $0 \leq j \lt q$, and write
\[
 A_{ij}(x)=\sum_{n\geq 0} a_{ij}(n)x^n,
\]
with $a_{ij}(n)=0$ for $n \lt 0$.
The infinite matrix $\operatorname{Lace}(A)$ has integer rows and columns.
If
\[
 u=pu' + i,\qquad v=qv' + j,\qquad 0\leq i \lt p,\quad 0\leq j \lt q,
\]
then
\[
 \operatorname{Lace}(A)_{u,v}=a_{ij}(v'-u').
\]
Equivalently, $\operatorname{Lace}(A)$ is the block Toeplitz
coefficient-convolution matrix associated to multiplication by $A$.
\end{definition}

For a column vector $A=(A_0,A_1,\dotsc,A_{p-1})^T$, Athanasiadis--Wagner call
the sequence \defin{fully interlacing} if $\operatorname{Lace}(A)$ is
\hyperref[totallyNonnegative]{totally non-negative}.  Thus every finite minor
of this infinite coefficient-convolution matrix is non-negative.

Full interlacing implies ordinary pairwise interlacing, but the converse is
false. The distinction is useful: pairwise interlacing is a root-level
condition, whereas full interlacing is a total-nonnegativity condition on a
structured coefficient matrix.


Whenever the product $AB$ is defined,
\[
 \operatorname{Lace}(AB)=\operatorname{Lace}(A)\operatorname{Lace}(B).
\]
This is convenient, since totally non-negative matrices are closed under multiplication.


\begin{example*}[Pairwise interlacing need not be fully interlacing]
The following example is from \cite[Ex.~3.4]{AthanasiadisWagner2024x}.  Let
\[
 P(x)=t+x,\qquad Q(x)=(b+x)(d+x),\qquad R(x)=(a+x)(c+x).
\]
If
\[
 a \leq b \leq t \leq c \leq d,
\]
then $P,Q,R$ are pairwise interlacing.  However, the $3\times 3$ minor
\[
 \begin{pmatrix}
 t   & 1   & 0 \\
 bd  & b+d & 1 \\
 ac  & a+c & 1
 \end{pmatrix}
\]
of $\operatorname{Lace}((P,Q,R)^T)$ has determinant
\[
 t(b+d-a-c)-(bd-ac),
\]
which can be negative.  For example, with
$a=1$, $b=2$, $t=2$, $c=3$, and $d=4$, this determinant is $-1$.
Thus ordinary pairwise interlacing does not imply full interlacing.
\end{example*}



\subsection[realMatrixRecursion]{Matrices preserving interlacing sequences}


The following very useful theorem is proved in \cite[Thm. 7.8.5]{Branden2015}.

\begin{theorem}[Matrices preserving interlacing sequences]
Let $G$ be an $m \times n$ matrix of polynomials in $t$.
Then $G : \mathscr{F}_n^+ \to \mathscr{F}_m^+$ if and only if
\begin{enumerate}
\item All entries of $G$ have nonnegative coefficients, and
\item For every $2{\times}2$ sub-matrix
$\begin{bmatrix}
 a(t) & b(t) \\
 c(t) & d(t)
\end{bmatrix}$ of $G$ and real numbers
$\lambda, \mu > 0$, we have
\[
 (\lambda t + \mu) b(t) + d(t) \interl (\lambda t + \mu) a(t) + c(t).
\]
\end{enumerate}
\end{theorem}


\begin{example*}[Minors preserving interleaving sequences]

There are 81 matrices with entries in $\{0,1,t\}$.
Remove those where first and last row are equal and those where one of the rows is $(0,0)$.
There are 56 such matrices left, and among these, the following 
satisfies the conditions in the theorem above.
\begin{equation*}
\begin{pmatrix}0 & 1 \\ 0 & t\end{pmatrix}, \quad
\begin{pmatrix}0 & 1 \\ t & 0\end{pmatrix}, \quad
\begin{pmatrix}1 & 0 \\ 0 & 1\end{pmatrix}, \quad
\begin{pmatrix}1 & 0 \\ t & 0\end{pmatrix}, \quad
\begin{pmatrix}t & 0 \\ 0 & t\end{pmatrix}, \quad
\begin{pmatrix}0 & 1 \\ t & 1\end{pmatrix}
\end{equation*}

\begin{equation*}
\begin{pmatrix}0 & 1 \\ t & t\end{pmatrix}, \quad
\begin{pmatrix}1 & 0 \\ 1 & 1\end{pmatrix}, \quad
\begin{pmatrix}1 & 0 \\ t & 1\end{pmatrix}, \quad
\begin{pmatrix}1 & 1 \\ 0 & 1\end{pmatrix}, \quad
\begin{pmatrix}1 & 1 \\ t & 0\end{pmatrix}, \quad
\begin{pmatrix}t & 0 \\ t & t\end{pmatrix}
\end{equation*}

\begin{equation*}
\begin{pmatrix}t & 1 \\ 0 & t\end{pmatrix}, \quad
\begin{pmatrix}t & 1 \\ t & 0\end{pmatrix}, \quad
\begin{pmatrix}t & t \\ 0 & t\end{pmatrix}, \quad
\begin{pmatrix}1 & 1 \\ t & 1\end{pmatrix}, \quad
\begin{pmatrix}1 & 1 \\ t & t\end{pmatrix}, \quad
\begin{pmatrix}t & 1 \\ t & t\end{pmatrix}
\end{equation*}
\end{example*}



\subsection[compatiblePolynomials]{Compatible polynomials}

Chudnovsky and Seymour introduced a concept closely related to interleaving sequences
of polynomials. The polynomials $f_1,\dotsc,f_k \in \setR[z]$ have a \defin{common interlacing}
if there is some $h \in \setR[z]$ with $h \interl f_j$ for all $j$.

\begin{theorem}[Obreschkoff--Dedieu, see
{\cite{Obreschkoff1963,Dedieu1992} and \cite[Thm.~8]{Branden2004operators}}]
Let $f,g \in \setR[x]$ have positive leading coefficients. Then every
polynomial in the real pencil
\[
  \{ \alpha f+\beta g : \alpha,\beta\in\setR\}
\]
is real-rooted or zero if and only if $f$ and $g$ are real-rooted and
one of $f \interl g$ or $g \interl f$ holds.
\end{theorem}

This signed pencil criterion is the two-polynomial counterpart of the
non-negative compatibility theorem below.

\begin{theorem}[Chudnovsky--Seymour, 2007, \cite{ChudnovskySeymour2007}]

Let $\{f_1,f_2,\dotsc,f_n\}$ be a set of real-rooted polynomials with \emph{positive} leading coefficients.
Such a set is said to be \defin{compatible} if for all
choices $\lambda_j\geq 0$, the sum $\sum_{j=1}^n \lambda_j f_j$ is real-rooted.

Then the following are equivalent:
\begin{itemize}
\item $f_1,\dotsc,f_n$ are pairwise compatible;
\item $f_1,\dotsc,f_n$ are compatible;
\item $f_i,f_j$ have a common interlacing for all $1\leq i \lt j \leq n$;
\item $f_1,\dotsc,f_n$ have a common interlacing.
\end{itemize}
\end{theorem}
This theorem can be generalized to polynomials with arbitrary real
coefficients (with negative leading coefficients), see \cite{LeakeRyder2024x}.

Another very interesting technique is presented in \cite{Bencs2014},
where Christoffel--Darboux type identities are used to show interlacing.
This exploits an idea similar to \hyperref[bezoutMatrix]{Bézout matrices}.
By taking the limit,
\[
\lim_{y\to x} \frac{f(x)g(y)-f(y)g(x)}{x-y} = f'(x)g(x) - g'(x)f(x)
\]
which is the Wronskian of $f$ and $g$.

\begin{proposition}[Wronskian criterion]
Let $f$ and $g$ be real-rooted polynomials with positive leading
coefficients. Then $g \interl f$ if and only if
\[
  f'(x)g(x)-f(x)g'(x)\geq 0
\]
for all $x\in\setR$.
\end{proposition}

Thus one way to prove an interlacing statement is to show that this
Wronskian is non-negative, perhaps by rewriting it as a sum of squares.
This is especially useful when a Christoffel--Darboux identity gives a
closed form for the difference quotient above.

\begin{theorem}[Interlacing families, see \cite{Lau2022InterlacingFamilies}]
Let $T$ be a rooted tree whose leaves are labelled by real-rooted
polynomials of the same degree and with positive leading coefficients.
Label every internal vertex by the sum of the polynomials labelling its
children. Suppose that the children of every internal vertex have a
common interlacing. Then, for each root index $j$, there are leaves
$a$ and $b$ such that
\[
  \lambda_j(p_a)\leq \lambda_j(p_{\emptyset})\leq \lambda_j(p_b),
\]
where $\lambda_j$ denotes the $j$th largest zero. In particular, some
leaf has largest zero at most the largest zero of the root polynomial.
\end{theorem}

For example, if $A$ is a real symmetric matrix and $A_S$ denotes the
principal submatrix obtained by deleting the rows and columns indexed by
$S$, then the characteristic polynomials $\{\det(xI-A_S): |S|=k\}$
form an interlacing family. The root polynomial is the $k$th derivative
of the characteristic polynomial of $A$. This refines the Cauchy
interlacing theorem from a single principal submatrix to a whole tree of
principal submatrices.


\section[interlacingTransforms]{Symmetric decompositions and preserving operators}


\subsection[symmetricDecompInterlacing]{Symmetric decompositions}

The following setup from \cite{BrandenSolus2018} is especially useful for
$h$-polynomials arising in \hyperref[ehrhart]{Ehrhart theory}.

Given $p \in \setR[x]$ with $\deg(p) \leq d$, define
\[
 I_d(p)(x) \coloneqq x^d p(1/x).
\]
Then there are unique polynomials $a,b \in \setR[x]$ such that
\[
 p(x) = a(x) + x b(x), \qquad I_d(a)=a, \qquad I_{d-1}(b)=b.
\]
This is the \defin{symmetric $I_d$-decomposition} of $p$, and it is given explicitly by
\[
 a = \frac{p - x I_d(p)}{1-x}, \qquad b = \frac{I_d(p)-p}{1-x}.
\]

It is also convenient to associate to $p$ the polynomial
\[
 f(x) \coloneqq (1+x)^d p\!\left( \frac{x}{1+x}\right),
\]
and to define
\[
 R_d(f)(x) \coloneqq (-1)^d f(-1-x).
\]
The paper uses the proper-position relation $\prec$; in the present non-negative coefficient
setting, this can be read as an interlacing relation.

\begin{theorem}[Brändén--Solus, \cite[Thm. 2.7]{BrandenSolus2018}]
Let $p \in \setR[x]$ have degree at most $d$, and let $(a,b)$ be its symmetric $I_d$-decomposition.
Assume that both $a$ and $b$ have non-negative coefficients. Then the following are equivalent:
\begin{itemize}
\item $b \prec a$;
\item $a \prec p$;
\item $b \prec p$;
\item $I_d(p) \prec p$;
\item $R_d(f) \prec f$.
\end{itemize}
In particular, these conditions imply that $p$, $a$, and $b$ are all real-rooted.
\end{theorem}

Thus, an interlacing statement between a polynomial and its reciprocal
forces the symmetric pieces of the decomposition to be real-rooted as well.
This is one reason the theorem shows up naturally for reflexive
$h^*$-polynomials; see \hyperref[magicPositiveEhrhart]{magic positivity} on the polytope page.



\subsection[realRootedPreservers]{Operators preserving real-rootedness}

A multivariate approach to preserving real-rootedness is provided by
\hyperref[stablePolynomial]{stable polynomials},
where \hyperref[stabilityOperators]{linear operators preserving stability}
imply real-rootedness of specializations.

\begin{theorem}[Hermite--Poulain, see
{\cite{Levin1964,Obreschkoff1963,Schur1914} and
\cite[Thm.~11]{Branden2004operators}}]
Let
\[
  f(x)=a_0+a_1x+\dotsb+a_nx^n
\]
and $g(x)$ be real-rooted polynomials. Then
\[
  f\!\left(\frac{d}{dx}\right)g
  \coloneqq
  a_0g(x)+a_1g'(x)+\dotsb+a_ng^{(n)}(x)
\]
is real-rooted.
\end{theorem}

In \cite{Wagner1992}, the following is proved.
\begin{theorem}
Suppose $f,g$ are polynomials, expressed as
\[
 \sum_{n \geq 0} f(n) z^n = \frac{W(f)(z)}{(1-z)^{1+\deg f}}, \quad
 \sum_{n \geq 0} g(n) z^n = \frac{W(g)(z)}{(1-z)^{1+\deg g}},
\]
with $W(f)(z)$ and $W(g)(z)$ being real-rooted with non-positive roots, then $W(fg)(z)$
also have real and non-positive roots.
\end{theorem}
This, in particular, shows that if $f$ and $g$ are \hyperref[ehrhart]{Ehrhart functions} of two polytopes $P_f$ and $P_g$
with real-rooted $h^*$-polynomials, then the $h^*$-polynomial of the polytope product $P_f * P_g$
is also real-rooted.
For examples of polytopes with real-rooted $h^*$-polynomials, see
\hyperref[realRootedTableaux]{tableaux and polytopes}.



\begin{theorem}[See \cite[Thm.~9]{Branden2004operators}]
If $\phi : \setR[x] \to \setR[x]$ is a linear operator that preserves real-rootedness,
then $\phi$ preserves the interleaving property:
\[
  f \interl g \implies \phi(f) \interl \phi(g).
\]
\end{theorem}





\subsection[derangedMap]{The deranged map}

Another operator with strong interlacing properties is the following.
Let $d_n(x)$ denote the classical derangement polynomial, and define the linear map
$D : \setR[x] \to \setR[x]$ by
\[
 D(x^n) \coloneqq d_n(x).
\]
If $p$ and $q$ are degree-$n$ polynomials with roots
$\alpha_1 \leq \dotsb \leq \alpha_n$ and $\beta_1 \leq \dotsb \leq \beta_n$
in $[-1,0]$, write $p \leq_{\mathrm{root}} q$ if $\alpha_k \leq \beta_k$ for all $k$.

\begin{theorem}[Brändén--Solus, \cite[Thm. 3.6]{BrandenSolus2018}]
If $p$ and $q$ have all roots in $[-1,0]$ and $p \leq_{\mathrm{root}} q$, then
$D(p) \prec D(q)$.
\end{theorem}

In particular, $D(p)$ and $D(q)$ are real-rooted and have interlacing zeros.
As a consequence, if
\[
 p(x)=\sum_{k=0}^n c_k x^k (x+1)^{n-k}, \qquad c_k \geq 0,
\]
then \cite[Cor. 3.7]{BrandenSolus2018} gives
\[
 A_n \prec D(p) \prec d_n.
\]
This is one route to the interlacing results for colored Eulerian and colored derangement
polynomials on \hyperref[realRootedWords]{the permutation page}.

There is an analogous basis-transformation viewpoint for other classical
polynomial triangles.  The Eulerian transformation \(A(x^n)=A_n(x)\) does not
preserve real-rootedness in complete generality \cite{BrandenJochemko2022},
but it does preserve real-rootedness on the non-negative cone discussed on the
\hyperref[gammaPositivity]{gamma-positivity page}
\cite{Athanasiadis2025EulerianTransformation}.  In parallel,
\name{Jianxi Mao} and \name{Lijie Wang} prove that the Narayana
transformation \(T_{N_m}(x^n)=N_{n,m}(x)\), interpolating between the type
\(B\) and type \(A\) Narayana polynomials, preserves real-rootedness for
polynomials with non-negative coefficients; see
\hyperref[narayanaTransformations]{Narayana transformations} and
\cite{MaoWang2026x}.

\name{Lily Li Liu} and \name{Xue Yan}
\cite{LiuYan2025InterlacingEulerianDerangement} refine this picture by
studying the triangular matrices arising from binomial-type transforms.
For the derangement transformation, if
\[
  D(t^k(1+t)^{n-k})=d_{n,k}(t),
\]
then, for every fixed $k$, the column
\[
  d_{k,k}(t),\,d_{k+1,k}(t),\,d_{k+2,k}(t),\dotsc
\]
is a generalized Sturm sequence, and for every fixed $n$ the row
\[
  d_{n,0}(t),\,d_{n,1}(t),\,\dotsc,\,d_{n,n}(t)
\]
is an interlacing sequence of real-rooted polynomials. They prove
analogous statements for the type $B$ and $r$-colored derangement
transformations, and obtain sufficient conditions for such row and
column interlacing in a more general recurrence setup.






\begin{remark*}[Pólya frequency sequences and TNN]
For the coefficient-level total-positivity background used above, see
\hyperref[polyaFrequency]{Pólya frequency sequences}.  In that language,
ordinary PF sequences are TNN Toeplitz matrices, while
\hyperref[fullyInterlacingLace]{fully interlacing sequences} are built from
block Toeplitz matrices.
\end{remark*}
